%% ---------------------------------------------------------------------------
%% There is no [[14,3,5]] quantum stabilizer code
%%
%% TODO before submission:
%%   1. Choose an arXiv license at submission time (see Section 9 notes).
%%   2. Optionally add a Zenodo DOI for the tagged artifact release.
%% ---------------------------------------------------------------------------
\documentclass[11pt]{amsart}

\usepackage{amssymb}
\usepackage[margin=1.2in]{geometry}
\usepackage{booktabs}
\usepackage[colorlinks=true,linkcolor=blue,citecolor=blue,urlcolor=blue]{hyperref}

\newtheorem{theorem}{Theorem}[section]
\newtheorem{lemma}[theorem]{Lemma}
\newtheorem{proposition}[theorem]{Proposition}
\newtheorem{corollary}[theorem]{Corollary}
\theoremstyle{definition}
\newtheorem{definition}[theorem]{Definition}
\theoremstyle{remark}
\newtheorem{remark}[theorem]{Remark}

\DeclareMathOperator{\wt}{wt}
\DeclareMathOperator{\rk}{rank}
\DeclareMathOperator{\Sh}{Sh}
\DeclareMathOperator{\supp}{supp}
\newcommand{\F}{\mathbb{F}}
\newcommand{\perps}{{\perp_s}}
\newcommand{\sform}[2]{\langle #1,#2\rangle_s}
\newcommand{\code}[1]{[\![#1]\!]}

\begin{document}

\title[There is no \code{14,3,5} code]{There is no $\code{14,3,5}$ quantum stabilizer code}

\author{Mansehej Singh}
\email{mansehej@yuzu.so}

\date{July 26, 2026}

\subjclass[2020]{81P73, 94B65, 94B05}
\keywords{quantum error correction, stabilizer codes, additive codes over GF(4),
shadow enumerators, nonexistence, computer-assisted proof}

\begin{abstract}
We prove that there is no binary quantum stabilizer code with parameters
$\code{14,3,5}$, degenerate codes included: no $11$-dimensional isotropic
subspace $C\subseteq\F_2^{28}$ satisfies
$\min\{\wt(v):v\in C^{\perps}\setminus C\}\ge 5$.  Consequently the optimal
minimum distance of a $14$-qubit stabilizer code encoding $3$ logical qubits
is exactly $4$, closing a gap that has stood in the code tables since
Calderbank, Rains, Shor, and Sloane (1998).  Since the linear-programming
bound for these parameters is $5$, the proof necessarily goes beyond linear
programming.  Candidates are split by the parity of stabilizer weights.
All-even stabilizers are eliminated by MacWilliams identities and
integrality alone.  For odd stabilizers we combine Rains' shadow enumerator
with two averaging identities over the $135$ self-dual extensions of the
code---for the weight enumerator, and, via an incidence computation in an
eight-dimensional orthogonal geometry of plus type, for the shadow---to
force one of three low-weight stabilizer profiles, each of which is then
eliminated by integrality or by counting coordinate collisions among the
weight-three shadow words.  Every target-specific computation is finite,
exact, and machine-verified by independent implementations; a one-command
replay artifact accompanies the paper.
\end{abstract}

\maketitle

\section{Introduction}\label{sec:intro}

A binary quantum stabilizer code $\code{n,k,d}$ encodes $k$ logical qubits
into $n$ physical qubits with minimum distance $d$
\cite{Gottesman1997,CRSS1998}.  For $n=14$ and $k=3$ the best construction
known has distance $4$, while the strongest upper bound known---the
linear-programming bound with shadow constraints---allows distance $5$.
The gap $4\le d_{\max}(14,3)\le 5$ appears already in Table~III of
Calderbank, Rains, Shor, and Sloane \cite{CRSS1998} and has remained open
in Grassl's online tables \cite{Grassl} since.  This paper closes it.

\begin{theorem}\label{thm:main}
There is no $11$-dimensional isotropic subspace $C\subseteq\F_2^{28}$ with
\[
\min\{\wt(v):v\in C^{\perps}\setminus C\}\ \ge\ 5 .
\]
Equivalently, there is no binary qubit stabilizer code $\code{14,3,5}$.
The statement includes degenerate (impure) codes and assumes no purity,
CSS structure, $\mathrm{GF}(4)$-linearity, or symmetry.
\end{theorem}

\begin{corollary}\label{cor:tables}
The optimal minimum distance of a binary stabilizer code with $n=14$ and
$k=3$ is exactly $4$, and the known $\code{14,3,4}$ codes are optimal.
\end{corollary}

Because the linear-programming bound admits a (rational) solution at
distance $5$, no proof of Theorem~\ref{thm:main} can proceed by linear
inequalities on weight enumerators alone; integrality and structural
information must enter.  Our argument is in the tradition of the
integrality refinements of \cite[\S 7]{CRSS1998} (e.g.\ the nonexistence
of a $\code{13,0,6}$ code) and of Rains' shadow theory
\cite{Rains1999Shadow}---compare the classical shadow bounds for self-dual
codes \cite{RainsSloane1998} and the nonexistence proof of Lam and Pless
\cite{LamPless1990}---and adds two ingredients that we believe are of
independent interest for small-parameter classification problems:

\begin{enumerate}
\item \emph{Self-dual extension averaging.}  A $\code{14,3}$ stabilizer,
viewed as an $11$-dimensional isotropic subspace $C$, lies in exactly $135$
self-dual (Lagrangian) extensions $D$, and the average of their weight
enumerators is the exact linear expression $(8A+B)/9$ in the enumerators
$A$ of $C$ and $B$ of $C^{\perps}$ (Lemma~\ref{lem:ext}).

\item \emph{Extension-shadow averaging.}  When the stabilizer contains
odd-weight elements, the shadows of those $135$ self-dual extensions
average, coefficientwise, to $\tfrac29 S_j$ for even $j$ and $0$ for odd
$j$, where $S$ is the shadow enumerator of $C$ itself
(Theorem~\ref{thm:tbar}).  The constant $\tfrac29=\tfrac{30}{135}$ is an
incidence number in the orthogonal geometry $O^+(8,2)$, verified by exhaustive
enumeration.
\end{enumerate}

Feeding both identities into the MacWilliams framework forces, for any
putative distance-$5$ code with an odd-weight stabilizer element, the
low-weight stabilizer profile
$(A_2,A_4)\in\{(0,1),(0,3),(1,1)\}$ with $A_1=A_3=0$
(Corollary~\ref{cor:profiles}); all-even stabilizers die earlier, by a
five-equation MacWilliams elimination (Proposition~\ref{prop:even}).  The
three surviving profiles are then eliminated in
Sections~\ref{sec:p11}--\ref{sec:p03}: profile $(1,1)$ forces the shadow
coefficient $S_2$ to be $\tfrac14$ or $\tfrac12$, contradicting
integrality; profiles $(0,1)$ and $(0,3)$ are destroyed by counting
\emph{collisions}---pairs of weight-three shadow words that agree in some
coordinate---which the shadow-coset structure simultaneously forces from
below and, after a local-Clifford normal-form classification of the
weight-four stabilizers, bounds from above.

The proof is computer-assisted in the following precise sense.  The
conceptual bridge---the lemmas of
Sections~\ref{sec:prelim}--\ref{sec:profiles}---is human-readable
mathematics and is proved in full in this paper.  Every target-specific
number (Krawtchouk coefficients, row reductions, the $135$ Lagrangians and
their incidences, the $O^+(8,2)$ shadow incidence, the $90$ translation
pairs, the $9{,}828$ weight-three words and their collision graphs) is a
finite exact computation, performed by five independently written
implementations in two languages, cross-compared, mutation-tested, and
replayable offline by one command (Section~\ref{sec:verification}).  No
floating-point arithmetic and no external solver are used anywhere.

\subsection*{Methodology disclosure}
The mathematical argument and all verification software were produced by
AI systems (Claude, Anthropic) in an autonomous research campaign, and were
then subjected to an extensive independent, adversarial, multi-agent
verification campaign (Section~\ref{sec:verification}).  In accordance with
arXiv policy no AI system is listed as an author; the human author has
reviewed the argument and assumes full responsibility for its correctness.

\section{Preliminaries}\label{sec:prelim}

\subsection{Stabilizer codes as isotropic subspaces}
Write vectors of $P=\F_2^{2n}$ as $v=(x\mid z)$ with $x,z\in\F_2^n$, and
equip $P$ with the symplectic form
\[
\sform{(x\mid z)}{(x'\mid z')}=x\cdot z'+z\cdot x' \pmod 2 .
\]
The \emph{Pauli weight} is
$\wt(x\mid z)=\#\{i:(x_i,z_i)\neq(0,0)\}$, via the coordinate map
$(0,0)\mapsto I$, $(1,0)\mapsto X$, $(0,1)\mapsto Z$, $(1,1)\mapsto Y$.
After discarding phases, an $n$-qubit stabilizer group with $r$ independent
generators and $-I$ not contained in it corresponds exactly to an
$r$-dimensional \emph{isotropic} subspace $C\subseteq P$ (two Paulis
commute iff their symplectic product vanishes), and conversely every
isotropic subspace lifts to such a group \cite{Gottesman1997,CRSS1998}.
(For a geometric survey of stabilizer codes see
\cite{BallCentellesHuber}.)  The code has $k=n-r$ logical qubits, and its
minimum distance is
\[
d=\min\{\wt(v):v\in C^{\perps}\setminus C\},
\]
where $C^{\perps}$ is the symplectic dual.  This is the standard
definition for possibly \emph{degenerate} codes: low-weight vectors
\emph{inside} $C$ are harmless, and only low-weight vectors of
$C^{\perps}\setminus C$ reduce the distance.  Throughout, $n=14$, $r=11$,
$k=3$, so $\dim C^{\perps}=17$, and we assume
\begin{equation}\label{eq:distance}
\min\{\wt(v):v\in C^{\perps}\setminus C\}\ \ge\ 5 ,
\end{equation}
deriving a contradiction.  Since a $\code{14,3,5}$ code satisfies
\eqref{eq:distance}, this proves Theorem~\ref{thm:main}.

\subsection{Weight enumerators and MacWilliams identities}
Let $A_i=\#\{c\in C:\wt(c)=i\}$ and $B_j=\#\{v\in C^{\perps}:\wt(v)=j\}$,
for $0\le i,j\le 14$.  Character orthogonality gives the MacWilliams
identities for additive codes \cite{CRSS1998}: with the quaternary
Krawtchouk polynomials
\begin{equation}\label{eq:kraw}
K_j(i)=\sum_{t}(-1)^t3^{\,j-t}\binom{i}{t}\binom{14-i}{j-t},
\end{equation}
we have
\begin{equation}\label{eq:mw}
B_j=\frac{1}{|C|}\sum_{i=0}^{14}A_iK_j(i),\qquad |C|=2^{11}=2048 .
\end{equation}
The inclusion $C\subseteq C^{\perps}$ gives $B_j\ge A_j$ for all $j$, and
the distance assumption \eqref{eq:distance} upgrades this to equality at
the bottom:
\begin{equation}\label{eq:dist-eq}
B_j=A_j \qquad (1\le j\le 4),
\end{equation}
because any word of $C^{\perps}$ of weight between $1$ and $4$ must lie in
$C$.  We emphasize that \eqref{eq:dist-eq} permits $A_1,\dots,A_4>0$
a priori: degeneracy is not excluded, it is carried through the argument.

\subsection{The parity quadratic form}
Define $q(v)=\wt(v)\bmod 2$.

\begin{lemma}\label{lem:q}
$q$ is a quadratic refinement of the symplectic form:
\[
q(u+v)=q(u)+q(v)+\sform{u}{v} \qquad (u,v\in P).
\]
Consequently $q$ restricted to an isotropic subspace $C$ is linear; if $C$
contains an odd-weight word, then $C_0=\ker(q|_C)$ has index $2$ in $C$,
and for our parameters $\dim C_0=10$.
\end{lemma}

\begin{proof}
Both sides are additive over coordinates, so it suffices to check the
identity on one coordinate, i.e.\ for the sixteen ordered pairs in
$\{I,X,Y,Z\}^2$: the weight parity of a product of two one-qubit Paulis
differs from the sum of their weight parities exactly when the two factors
are distinct non-identity Paulis, which is exactly when they anticommute.
(The artifact check \texttt{ALG.QFORM.POLAR} repeats all sixteen cases.)
On isotropic $C$ the polar term vanishes, so $q|_C$ is linear.
\end{proof}

Under the orthogonal sum decomposition of $(P,q)$ into $14$ coordinate
planes, each plane is anisotropic (all three nonzero vectors have $q=1$),
hence has Arf invariant $1$; the Arf invariant of $P$ is therefore
$14\bmod 2=0$, i.e.\ $(P,q)$ is of \emph{plus type}.  We use this in
Lemma~\ref{lem:witt}.

Every candidate $C$ now falls into exactly one of two branches: the
\emph{all-even branch} ($q|_C\equiv 0$) and the \emph{odd branch}
($C$ contains an odd-weight word).  Section~\ref{sec:even} eliminates the
former; the rest of the paper eliminates the latter.

\section{The all-even branch}\label{sec:even}

\begin{proposition}\label{prop:even}
No $11$-dimensional isotropic $C\subseteq\F_2^{28}$ with all weights even
satisfies \eqref{eq:distance}.
\end{proposition}

\begin{proof}
Assume the contrary.  Then $A_i=0$ for odd $i$, leaving the seven unknowns
$A_2,A_4,\dots,A_{14}$ (with $A_0=1$), subject to five affine equations:
$\sum_{i}A_i=2048$, and the four distance equations obtained by combining
\eqref{eq:mw} with \eqref{eq:dist-eq},
\[
\sum_{i\ \mathrm{even}}A_i\bigl(K_j(i)-2048\,\delta_{ij}\bigr)=-K_j(0)
\qquad (1\le j\le 4).
\]
Exact rational row reduction of this system (artifact check
\texttt{EVEN.RREF}) leaves the two free parameters $p=A_{12}$, $q=A_{14}$
and expresses
\[
40A_2=-1145+p+6q,\qquad
10A_4=-35+p-14q,\qquad
20A_6=4965-13p+102q .
\]
Multiplying the target combination by ten and substituting,
\[
160A_2+90A_4+20A_6
=4(-1145+p+6q)+9(-35+p-14q)+(4965-13p+102q)=70 ,
\]
where both $p$ and $q$ cancel.  Hence every solution satisfies the single
relation
\begin{equation}\label{eq:even-relation}
16A_2+9A_4+2A_6=7 .
\end{equation}
All coefficients are positive and all $A_i$ are nonnegative integers, so
$16A_2\le 7$ forces $A_2=0$, then $9A_4\le7$ forces $A_4=0$, and
$2A_6=7$ is impossible modulo $2$.
\end{proof}

\section{Shadows, self-dual extensions, and the averaging identities}
\label{sec:shadow}

For the remainder of the paper $C$ is an $11$-dimensional isotropic
subspace satisfying \eqref{eq:distance} whose weight function is not
identically even, and $C_0=\ker(q|_C)$, $\dim C_0=10$.

\subsection{The shadow of $C$}

\begin{definition}
The \emph{shadow} of $C$ is the set
$\mathcal S=C_0^{\perps}\setminus C^{\perps}$, with weight enumerator
$S_j=\#\{s\in\mathcal S:\wt(s)=j\}$.
\end{definition}

\begin{lemma}\label{lem:shadow}
Let $E$ be the weight enumerator of $C_0^{\perps}$.  Then
\begin{equation}\label{eq:eshadow}
E_j=\frac{1}{|C_0|}\sum_{i\ \mathrm{even}}A_iK_j(i),
\qquad
S_j=E_j-B_j=\frac{1}{|C|}\sum_{i=0}^{14}(-1)^iA_iK_j(i) .
\end{equation}
Moreover $\mathcal S$ is the nontrivial coset of $C^{\perps}$ in
$C_0^{\perps}$; in particular
$\mathcal S+\mathcal S\subseteq C^{\perps}$ and
$\mathcal S+C^{\perps}=\mathcal S$, every shadow word commutes with
$C_0$, and every shadow word has symplectic product $1$ with every
odd-weight word of $C$.
\end{lemma}

\begin{proof}
The first formula is \eqref{eq:mw} applied to $C_0$, whose weights are
exactly the even weights of $C$ with the same multiplicities, and
$|C_0|=|C|/2$.  Since $C_0\subset C$ has index $2$, dually
$C^{\perps}\subset C_0^{\perps}$ has index $2$; the shadow is the
nontrivial coset and $S=E-B$; subtracting \eqref{eq:mw} from
\eqref{eq:eshadow} flips the sign of the odd-$i$ terms, giving the signed
transform.  The coset properties
are immediate, and a shadow word $s$ lies in $C_0^{\perps}$ but not
$C^{\perps}$, so $\sform{s}{\cdot}$ vanishes on $C_0$ and is $1$ on the odd
coset $C\setminus C_0$.
\end{proof}

This is Rains' quantum shadow \cite{Rains1999Shadow} in binary symplectic
language.  Note $S_j\ge0$, and $E_j=B_j+S_j\ge A_j$, always.

\subsection{Self-dual extensions and the enumerator average}

The quotient $V=C^{\perps}/C$ is a nondegenerate symplectic space of
dimension $2k=6$.

\begin{lemma}\label{lem:ext}
Self-dual subspaces $D$ with $C\subset D=D^{\perps}$ correspond bijectively
to Lagrangian (maximal isotropic, $3$-dimensional) subspaces of $V$.
There are exactly $135$ of them; each nonzero point of $V$ lies in exactly
$15$.  Consequently, writing $N^{(D)}$ for the weight enumerator of $D$,
\begin{equation}\label{eq:nbar}
\overline N \;=\;\frac{1}{135}\sum_{D}N^{(D)}\;=\;\frac{8A+B}{9}
\qquad\text{coefficientwise.}
\end{equation}
\end{lemma}

\begin{proof}
The symplectic form descends to $V$ and is nondegenerate.  If
$\bar L\subseteq V$ is Lagrangian, its preimage $D$ satisfies
$\dim D=11+3=14=n$ and is isotropic, hence $D=D^{\perps}$; the
correspondence inverts by taking quotients.  The counts---$135$
Lagrangians among the $1395$ three-dimensional subspaces of $\F_2^6$, with
every one of the $63$ nonzero points on exactly $15$ of them---are an
exhaustive enumeration (artifact check \texttt{EXT.LAGRANGIAN\_6D};
$135\cdot 7=63\cdot 15=945$).  For the average: a word of $C$ lies in all
$135$ extensions, while a word of $C^{\perps}\setminus C$ maps to a
nonzero point of $V$ and lies in exactly the $15$ extensions through that
point.  Hence
$\sum_D N^{(D)}=135A+15(B-A)$, which is \eqref{eq:nbar}.
\end{proof}

Each $D$ is a self-dual additive code of length $14$
\cite{CRSS1998,NRS2006}; its enumerator is a fixed point of the
MacWilliams transform.  Applying the transform to
\eqref{eq:nbar} yields an identity that is automatically satisfied, so
\eqref{eq:nbar} alone adds no linear information beyond \eqref{eq:mw}.
The new information comes from the \emph{shadows} of the extensions.

\subsection{The quotient geometry and the extension--shadow average}

Every extension $D$ contains the odd coset of $C$, hence is an odd code;
let $D_0$ be its even kernel ($\dim D_0=13$) and
$\Sh(D)=D_0^{\perps}\setminus D$ its shadow, a set of $2^{14}$ words with
enumerator $T^{(D)}$.  By Lemma~\ref{lem:shadow} applied to $D$,
\begin{equation}\label{eq:tshadow}
T^{(D)}_j=\frac{1}{2^{14}}\sum_i(-1)^iN^{(D)}_iK_j(i).
\end{equation}

Consider now $W=C_0^{\perps}/C_0$, a nondegenerate symplectic space of
dimension $8$.  Since $q(v+c)=q(v)+q(c)+\sform{v}{c}=q(v)$ for $c\in C_0$
and $v\in C_0^{\perps}$, the parity form $q$ descends to $W$.

\begin{lemma}[Witt reduction]\label{lem:witt}
$(W,q)$ is a nondegenerate quadratic space of plus type, isometric to
$\F_2^4\oplus\F_2^4$ with $q(x\mid z)=x\cdot z$.  The class
$r=C\setminus C_0\in W$ is anisotropic: $q(r)=1$.
\end{lemma}

\begin{proof}
Choose a nonzero $e\in C_0$; it is singular ($q(e)=0$).  By nondegeneracy
of $P$ pick $g$ with $\sform{e}{g}=1$; if $q(g)=1$ replace $g$ by $g+e$,
which satisfies $q(g+e)=q(g)+q(e)+\sform{g}{e}=0$.  For every other basis
vector $c$ of $C_0$ with $\sform{c}{g}=1$, replace $c$ by $c+e$: this
keeps it singular ($q(c+e)=q(c)+q(e)+\sform{c}{e}=0$) and makes it
perpendicular to the hyperbolic plane $H=\langle e,g\rangle$.  Restrict to
$H^{\perp}$ (again nondegenerate) and iterate through a basis of $C_0$.
After ten steps we have split off ten hyperbolic planes, each of Arf
invariant $0$, containing a basis of $C_0$; the Arf invariant is additive
over orthogonal sums and classifies nondegenerate quadratic forms over
$\F_2$ \cite{Grove2002}, and the ambient Arf invariant is $0$
(Section~\ref{sec:prelim}), so the orthogonal remainder $R$ of dimension
$28-20=8$ has Arf invariant $0$, i.e.\ is of plus type.  Finally
$C_0^{\perps}=C_0\oplus R$, so $W\cong R$ as quadratic spaces.  The class
$r$ consists of odd-weight words, so $q(r)=1$.
\end{proof}

Under the correspondence of Lemma~\ref{lem:ext}, an extension $D$
determines the Lagrangian $\bar D=D/C_0$ of $W$ (dimension $4$), which
contains $r$; this is again a bijection onto the $135$ Lagrangians of $W$
through $r$.

\begin{lemma}\label{lem:extshadow}
For each extension $D$, the shadow $\Sh(D)$ is a union of $C_0$-cosets,
and its image in $W$ is the affine solution set
\[
\Sh(\bar D)=\{h\in W:\sform{h}{\bar d}=q(\bar d)\ \text{for all }
\bar d\in\bar D\},
\]
a coset of $\bar D$ of size $16$.  Exhaustive enumeration in the canonical
model of Lemma~\ref{lem:witt} (artifact check
\texttt{EXT.SHADOW\_INCIDENCE\_8D}) shows: every class in every
$\Sh(\bar D)$ is singular ($q=0$); among the $128$ classes $h$ with
$\sform{h}{r}=1$, exactly $72$ are singular, and each of these lies in
exactly $30$ of the $135$ sets $\Sh(\bar D)$, while the $56$ nonsingular
ones lie in none ($135\cdot16=72\cdot30=2160$).  Moreover the orthogonal
group of $(W,q)$ acts transitively on the $120$ anisotropic vectors (the
transvection orbit of $r$ is enumerated), so the count is independent of
the choice of $r$ and transfers to the actual quotient by
Lemma~\ref{lem:witt}.
\end{lemma}

\begin{proof}
$C_0\subseteq D_0$ and $C_0\subseteq D\subseteq D_0^{\perps}$, so both
$D_0^{\perps}$ and $D$ are unions of $C_0$-cosets, hence so is $\Sh(D)$.
For the affine description, work in $W$: $\bar D$ is Lagrangian and
$\bar D_0=\ker(q|_{\bar D})$ has dimension $3$, so
$\Sh(D)/C_0=\bar D_0^{\perp}\setminus\bar D$, of size $2^5-2^4=16$.  A
class $h$ lies in $\bar D_0^{\perp}$ iff $\sform{h}{\bar d}=0=q(\bar d)$
for all even $\bar d\in\bar D$, and then $h\notin\bar D=\bar D^{\perp}$
iff $\sform{h}{\bar d}=1=q(\bar d)$ on the odd coset of $\bar D$.  The
solution set of the affine system $\sform{h}{\cdot}=q(\cdot)$ on $\bar D$
is a coset of $\bar D^{\perp}=\bar D$.  The incidence counts and the
transvection orbit are finite exhaustive computations; transitivity on
anisotropic vectors supplies an isometry carrying the actual $(W,q,r)$ to
the canonical model, under which the incidence numbers are preserved.
\end{proof}

\begin{theorem}[Extension--shadow average]\label{thm:tbar}
Let $\overline T=\frac{1}{135}\sum_D T^{(D)}$.  Then
\begin{equation}\label{eq:tbar}
\overline T_j=
\begin{cases}
\dfrac{2}{9}\,S_j, & j\ \text{even},\\[4pt]
0, & j\ \text{odd}.
\end{cases}
\end{equation}
\end{theorem}

\begin{proof}
Fix $j$ and count the pairs $(D,s)$ with $s\in\Sh(D)$, $\wt(s)=j$.  Every
word of every $\Sh(D)$ lies in $C_0^{\perps}$ (as
$D_0^{\perps}\subseteq C_0^{\perps}$) and has
$\sform{s}{r}=q(r)=1$ by the affine description, hence lies in the shadow
$\mathcal S$ of $C$.  Conversely a shadow word $s$ of weight $j$ has class
$\bar s$ with $\sform{\bar s}{r}=1$ (Lemma~\ref{lem:shadow}) and
$q(\bar s)=j\bmod 2$.  If $j$ is even, Lemma~\ref{lem:extshadow} says
$\bar s$ lies in exactly $30$ of the sets $\Sh(\bar D)$, so $s$ lies in
exactly $30$ of the sets $\Sh(D)$; if $j$ is odd, in none.  Hence
$\sum_DT^{(D)}_j=30\,S_j$ for even $j$ and $0$ for odd $j$; divide by
$135$.
\end{proof}

Combining \eqref{eq:tshadow}, \eqref{eq:nbar}, and \eqref{eq:tbar} yields,
for every $j$, an exact linear equation in the $A_i$:
\begin{equation}\label{eq:family}
\frac{1}{2^{14}}\sum_i(-1)^i\Bigl(\frac{8A+B}{9}\Bigr)_iK_j(i)
=\begin{cases}\frac{2}{9}S_j,& j\ \text{even},\\ 0,& j\ \text{odd},\end{cases}
\qquad B,S\ \text{as in \eqref{eq:mw}, \eqref{eq:eshadow}}.
\end{equation}
These fifteen equations, together with \eqref{eq:mw},
\eqref{eq:dist-eq}, $A_0=1$, $\sum_iA_i=2048$, and
$\sum_i(-1)^iA_i=0$ (the odd branch has $|C_0|=|C|/2$), constitute the
\emph{odd-branch constraint family} used below.

\section{The coefficient identity and the three profiles}
\label{sec:profiles}

\begin{proposition}[Identity 51]\label{prop:id51}
Every $C$ in the odd branch satisfying \eqref{eq:distance} obeys
\begin{equation}\label{eq:id51}
168A_1+28A_2+63A_3+15A_4+14A_5+4A_6+2A_{14}+4E_1+2E_4=51 ,
\end{equation}
where $E=B+S$ is the enumerator of $C_0^{\perps}$ as in
Lemma~\ref{lem:shadow}.
\end{proposition}

\begin{proof}
This is an exact linear consequence of the elementary subfamily alone:
$\sum_{i\ge1}A_i=2047$, $\sum_{i\ \mathrm{even},\ge2}A_i=1023$, the four
distance rows $\sum_{i\ge1}(K_j(i)-2048\delta_{ij})A_i=-K_j(0)$
($1\le j\le4$), and the two defining rows
$\sum_{i\ \mathrm{even},\ge2}K_j(i)A_i-1024E_j=-K_j(0)$ ($j\in\{1,4\}$)
from \eqref{eq:eshadow}.  The eliminating multipliers are
\[
(-147,\,-148,\,-42,\,-42,\,-7,\,-7,\,8,\,4),
\]
in the order listed, followed by division by $-2048$: expanding this
combination with the Krawtchouk values \eqref{eq:kraw} produces
\eqref{eq:id51} exactly.  The expansion is a mechanical rational
computation, re-derived from \eqref{eq:kraw} at every artifact run
(check \texttt{ODD.IDENTITY51}) and independently by the further
implementations of Section~\ref{sec:verification}.
\end{proof}

\begin{corollary}[Profile trichotomy]\label{cor:profiles}
In the odd branch, $A_1=A_3=0$ and
\[
(A_2,A_4)\in\{(0,1),\,(0,3),\,(1,1)\} .
\]
\end{corollary}

\begin{proof}
All quantities in \eqref{eq:id51} are nonnegative integers.  From
$168A_1\le 51$ and $63A_3\le 51$ we get $A_1=A_3=0$.  The left side is
even except for the term $15A_4$, and the right side is odd, so $A_4$ is
odd; $15\cdot 5>51$ leaves $A_4\in\{1,3\}$.  If $A_4=3$ the remaining
budget is $51-45=6$, so $28A_2\le6$ gives $A_2=0$.  If $A_4=1$ the budget
is $36$ and $28A_2\le 36$ gives $A_2\in\{0,1\}$.
\end{proof}

The full odd-branch family (including \eqref{eq:family}) then pins the two
remaining top coefficients.  The following is again an exact elimination,
machine-verified as an implication (not merely a consistency check) by
row reduction of the family (artifact check
\texttt{ANTI.UNIVERSAL\_BRANCH\_FORMULAS}).

\begin{proposition}[Universal branch formulas]\label{prop:universal}
In the odd branch (so $A_1=A_3=0$),
\begin{align}
A_{14}&=36-18A_2-12A_4-12A_5-3A_6-4S_2-2S_4 , \label{eq:a14}\\
S_3&=\tfrac{105}{2}-6A_2-\tfrac{19}{2}A_4-19A_5-3A_6-18S_2-4S_4 .
\label{eq:s3}
\end{align}
\end{proposition}

All variables appearing here are nonnegative integers (they count words).
The three profiles are now eliminated in turn.

\section{Profile $(1,1)$}\label{sec:p11}

\begin{proposition}\label{prop:p11}
Profile $(A_2,A_4)=(1,1)$ is impossible.
\end{proposition}

\begin{proof}
Let $w_2$ and $w_4$ be the unique stabilizer words of weights $2$ and $4$.
Both are even, lie in $C_0$, and commute; writing $a$ for the number of
coordinates where they agree (with equal nonidentity Pauli) and $b$ for
the number where both are nonidentity but different, commutation forces
$b$ even, and $\wt(w_2+w_4)=6-2a-b\in\{2,4,6\}$.  Weight $2$ would force
$w_2+w_4=w_2$ (as $A_1=0$, $A_2=1$), i.e.\ $w_4=0$; weight $4$ would force
$w_2+w_4=w_4$, i.e.\ $w_2=0$.  Hence $\wt(w_2+w_4)=6$, the supports are
disjoint, and $A_6\ge1$.

Substituting $A_2=A_4=1$, $E_1=S_1$, $E_4=1+S_4$ into \eqref{eq:id51}
gives
\[
14A_5+4A_6+2A_{14}+4S_1+2S_4=6 .
\]
With $A_6\ge1$: $A_5=0$, $A_6=1$, $S_1=0$, and
$(A_{14},S_4)\in\{(0,1),(1,0)\}$.  Now \eqref{eq:a14} with
$A_2=A_4=A_6=1$, $A_5=0$ reads $A_{14}=3-4S_2-2S_4$.  The case
$(A_{14},S_4)=(0,1)$ forces $S_2=\tfrac14$; the case $(1,0)$ forces
$S_2=\tfrac12$.  Both contradict the integrality of the shadow coefficient
$S_2$ (artifact check \texttt{P11.INTEGRALITY}).
\end{proof}

\section{Profile $(0,1)$}\label{sec:p01}

Here $A_1=A_2=A_3=0$ and $A_4=1$; write $u$ for the unique weight-four
stabilizer word.  Substituting into \eqref{eq:id51},
\eqref{eq:a14}, \eqref{eq:s3} and eliminating $A_6$ gives the exact
formulas (artifact check \texttt{P01.FORMULAS})
\begin{align}
A_6&=7-5A_5-4S_2-S_4+2S_1, \label{eq:p01-a6}\\
A_{14}&=3+3A_5+8S_2+S_4-6S_1, \label{eq:p01-a14}\\
S_3&=22-4A_5-6S_2-S_4-6S_1. \label{eq:p01-s3}
\end{align}

We now bound the shadow from the combinatorial side.  Call a
\emph{bin} a pair (coordinate, nonidentity Pauli); there are
$14\cdot3=42$ bins, and a weight-three word occupies $3$ bins.

\begin{lemma}\label{lem:p01-s1}
$S_1\le 1$.  If $S_1=1$ then $S_2=0$ and $S_3\le1$.
\end{lemma}

\begin{proof}
Two distinct weight-one shadow words differ by a word of $C^{\perps}$
(Lemma~\ref{lem:shadow}) of weight $\le2$, hence by
\eqref{eq:distance} a nonzero word of $C$ of weight $\le2$, contradicting
$A_1=A_2=0$.  If $h$ has weight $1$ and $h'$ weight $2$, both shadow
words, then $h+h'\in C$ has weight between $1$ and $3$, contradicting
$A_1=A_2=A_3=0$; so $S_2=0$.  For weight-three shadow $g$, the word
$h+g\in C$ has weight $\le4$, hence $h+g=u$ and $g=h+u$ is unique.
\end{proof}

\begin{lemma}\label{lem:p01-bins}
Suppose $S_1=0$.  Then among the weight-three shadow words at most one
pair occupies a common bin, and consequently
\[
3S_3\le 42+1,\quad\text{i.e.}\quad S_3\le 14 .
\]
If moreover $S_2=1$, then every weight-three shadow word has support
disjoint from the weight-two shadow word, only $36$ bins are available,
and $3S_3\le 36+1$, i.e.\ $S_3\le 12$.
\end{lemma}

\begin{proof}
If two distinct weight-three shadow words $g,g'$ share a bin, then
$\wt(g+g')\le 3+3-2=4$ and $g+g'\in C\setminus\{0\}$, so $g+g'=u$: shared-bin
pairs are \emph{translation pairs} $\{g,g+u\}$.  Normalize
$u=X_1X_2X_3X_4$ by a coordinate permutation and one-qubit Cliffords
(these preserve weights, the symplectic form, membership in $C$, $C_0$,
$C^{\perps}$, and hence every hypothesis; see also
Lemma~\ref{lem:lc}).  Exhaustive enumeration (artifact checks
\texttt{P01.COLLISION\_90}, \texttt{P01.S1\_TRANSLATIONS}) shows: there
are exactly $90$ translation pairs with both members of weight three
sharing a bin; each shares exactly one bin; and among the
$\binom{90}{2}=4005$ unordered pairs of distinct translation pairs,
\emph{none} is compatible, where compatibility means all four cross
differences between the two pairs are either $u$ or of weight $\ge5$---the
condition forced on shadow words by \eqref{eq:distance} and
$A_1=A_2=A_3=0$, $A_4=1$.  Hence at most one translation pair occurs among
the shadow words, so at most one bin is occupied twice and none thrice
(a thrice-occupied bin would contain two distinct translation pairs).
Counting incidences, $3S_3\le 42+1$.

For the second claim let $h$ be the weight-two shadow word and $g$ any
weight-three shadow word.  Then $h+g\in C^{\perps}$ has weight $\le5$; if
the weight were $\le4$, then $h+g$ would be a nonzero word of $C$, hence
equal to $u$, giving $g=h+u$.  But $u\in C_0$ and $h\in C_0^{\perps}$
give $\sform{h}{u}=0$ and $q(u)=0$, so by Lemma~\ref{lem:q}
$q(h+u)=q(h)=0$, whereas $q(g)=1$ since $\wt(g)=3$---a contradiction.  So
$\wt(h+g)=5$ exactly, which for weights $2$ and $3$ means
$\supp(h)\cap\supp(g)=\emptyset$ (this equivalence is also verified
exhaustively over all $9828$ weight-three words, artifact check
\texttt{P01.BINS\_36}).  All weight-three shadow words therefore live on
the twelve coordinates avoiding $\supp(h)$: $36$ bins, and the repeated-bin
argument gives $3S_3\le37$.
\end{proof}

\begin{proposition}\label{prop:p01}
Profile $(A_2,A_4)=(0,1)$ is impossible.
\end{proposition}

\begin{proof}
All variables below are nonnegative integers.  By Lemma~\ref{lem:p01-s1},
$S_1\le1$ and $S_1=1$ forces $S_2=0$; and when $S_1=0$,
\eqref{eq:p01-a6} gives $4S_2\le7$, so $S_2\le1$.  Exactly three cases
remain.

\emph{Case $S_1=1$ (so $S_2=0$, $S_3\le1$).}  By \eqref{eq:p01-s3},
$S_3=16-4A_5-S_4\le 1$ forces $4A_5+S_4\ge15$, while \eqref{eq:p01-a6}
gives $A_6=9-5A_5-S_4\ge0$, i.e.\ $5A_5+S_4\le9$.  Then $A_5\le1$; for
$A_5=0$ we would need $S_4\ge15$ against $S_4\le9$, and for $A_5=1$ we
would need $S_4\ge11$ against $S_4\le4$.  Empty.

\emph{Case $S_1=0$, $S_2=0$ (so $S_3\le14$).}  Now $S_3=22-4A_5-S_4\le14$
forces $4A_5+S_4\ge8$, while $A_6=7-5A_5-S_4\ge0$ forces $5A_5+S_4\le7$:
for $A_5=0$, $S_4\ge8$ versus $S_4\le7$; for $A_5=1$, $S_4\ge4$ versus
$S_4\le2$; $A_5\ge2$ is excluded.  Empty.

\emph{Case $S_1=0$, $S_2=1$ (so $S_3\le12$).}  Now
$S_3=16-4A_5-S_4\le12$ forces $4A_5+S_4\ge4$, while
$A_6=3-5A_5-S_4\ge0$ forces $A_5=0$ and $S_4\le3$, whence $S_4\ge4$,
a contradiction.  Empty.

The three cases are exhaustive, so no code with profile $(0,1)$ exists.
(The same ledger is exhausted mechanically, over ranges wider than needed,
in artifact check \texttt{P01.EXHAUSTED}; removing either bin bound of
Lemma~\ref{lem:p01-bins} leaves algebraic survivors, so both are genuinely
load-bearing.)
\end{proof}

\section{Profile $(0,3)$}\label{sec:p03}

Here $A_1=A_2=A_3=0$ and $A_4=3$: there are exactly three weight-four
stabilizer words $u,v,w$, all in $C_0$, pairwise commuting.

\begin{lemma}\label{lem:p03-forced}
In profile $(0,3)$,
\[
A_5=A_6=A_{14}=S_1=S_2=S_4=0
\qquad\text{and}\qquad
S_3=24 .
\]
\end{lemma}

\begin{proof}
Substituting $A_2=0$, $A_4=3$, $E_1=S_1$, $E_4=3+S_4$ into
\eqref{eq:id51} gives $45+6+14A_5+4A_6+2A_{14}+4S_1+2S_4=51$, i.e.
\[
14A_5+4A_6+2A_{14}+4S_1+2S_4=0 ,
\]
so $A_5=A_6=A_{14}=S_1=S_4=0$.  Then \eqref{eq:a14} reads
$0=36-36-4S_2$, forcing $S_2=0$, and \eqref{eq:s3} gives
$S_3=\tfrac{105}{2}-\tfrac{57}{2}=24$.
\end{proof}

\begin{lemma}[Local Clifford normal forms]\label{lem:lc}
Coordinate permutations and independent one-qubit Clifford operations act
on $P$ preserving weight, the symplectic form, and $q$; on each coordinate
the one-qubit Clifford group induces the full symmetric group $S_3$ on
$\{X,Y,Z\}$.  Up to these operations, the triple $\{u,v,w\}$ is one of
exactly four forms:
\begin{itemize}
\item if $\rk\{u,v,w\}=2$ (i.e.\ $w=u+v$): one of
\[
\mathrm{p0}=\begin{pmatrix}XXXX\\ZZZZ\\YYYY\end{pmatrix},\quad
\mathrm{p1}=\begin{pmatrix}XIXXX\\IXXZZ\\XXIYY\end{pmatrix},\quad
\mathrm{p2}=\begin{pmatrix}XXIIXX\\IIXXXX\\XXXXII\end{pmatrix}
\]
(identity on all unshown coordinates);
\item if $\rk\{u,v,w\}=3$: the three supports are pairwise disjoint and
the triple is equivalent to $X^{\otimes4}$ on coordinates $1$--$4$,
$5$--$8$, $9$--$12$.
\end{itemize}
\end{lemma}

\begin{proof}
In the rank-two case classify the coordinates in
$\supp(u)\cup\supp(v)$ by the (unordered up to $S_3$-relabelling) type of
the column $(u_i,v_i)$: the possibilities are $(a,I)$, $(I,b)$, $(a,a)$,
$(a,b)$ with $a\ne b$ nonidentity, occurring $A,B,C,D$ times.  The three
nonzero words $u$, $v$, $u+v$ have weights $A+C+D$, $B+C+D$, $A+B+D$, all
equal to $4$, whence $A=B=C=p$ and $D=4-2p$ with $p\in\{0,1,2\}$
(commutation is automatic: $\sform{u}{v}=D\bmod2=0$).  Per-coordinate
$S_3$ freedom maps each column of its type to a fixed representative and
a coordinate permutation sorts the columns, producing exactly
$\mathrm{p0},\mathrm{p1},\mathrm{p2}$ for $p=0,1,2$.  In the rank-three
case, for any two of the three words the sum $u+v\in C_0$ is nonzero of
even weight; weight $\le4$ would put it in $\{u,v,w\}$ (as
$A_1=A_2=A_3=0$, $A_4=3$), forcing rank two; weight $6$ is excluded by
$A_6=0$ (Lemma~\ref{lem:p03-forced}); so $\wt(u+v)=8$ and the supports
are disjoint, and each column is normalized to $X$.  Exhaustiveness of
this classification is confirmed independently by brute force: fixing
$u=X_1X_2X_3X_4$, all $81{,}081$ weight-four words $v$ are enumerated and
canonicalized (artifact checks \texttt{P03.NORMAL\_FORMS\_4} and the
JavaScript clean-room enumeration).
\end{proof}

\begin{proposition}\label{prop:p03}
Profile $(A_2,A_4)=(0,3)$ is impossible; with
Propositions~\ref{prop:even}, \ref{prop:p11}, and \ref{prop:p01} this
completes the proof of Theorem~\ref{thm:main}.
\end{proposition}

\begin{proof}
By Lemma~\ref{lem:p03-forced} there are exactly $24$ weight-three shadow
words.  Each occupies $3$ of the $42$ bins, giving $72$ incidences, so the
number of unordered pairs of shadow words sharing a bin is at least
\[
\sum_{\text{bins}}\binom{n_b}{2}\ \ge\ \sum_{\text{bins}}(n_b-1)\ \ge\
72-42=30 .
\]
As in Lemma~\ref{lem:p01-bins}, any two distinct weight-three shadow words
sharing a bin differ by a nonzero word of $C$ of weight $\le4$, i.e.\ by
one of $u,v,w$.  Now fix one of the four normal forms of
Lemma~\ref{lem:lc} and consider the graph whose vertices are the
weight-three words commuting with $u,v,w$ (shadow words commute with
$C_0\supseteq\{u,v,w\}$, so the true shadow words are among these) and
whose edges join $g$ and $g+x$ ($x\in\{u,v,w\}$) when both have weight
three and share a bin.  Exhaustive enumeration of all $9{,}828$
weight-three words (artifact checks \texttt{P03.WORDS\_9828},
\texttt{P03.DEGREE\_HISTOGRAMS}, \texttt{P03.COLLISION\_CONTRADICTION},
independently in three implementations) yields, for the four forms, vertex
degree histograms whose top twenty-four degrees sum to at most
\[
24,\quad 24,\quad 26,\quad 56 ,
\]
so that \emph{any} $24$ vertices induce at most
\[
12,\quad 12,\quad 13,\quad 28
\]
edges, respectively.  Every shared-bin pair among the $24$ shadow words is
an induced edge, so there are at most $28<30$ such pairs.  This
contradicts the lower bound above in every normal form, and profile
$(0,3)$ is impossible.
\end{proof}

\section{Computer verification}\label{sec:verification}

\subsection{What is trusted and what is machine-checked}
The handwritten kernel of the proof consists of the lemmas proved in full
above: the stabilizer correspondence, Lemma~\ref{lem:q},
Proposition~\ref{prop:even} (given its displayed row reduction),
Lemma~\ref{lem:shadow}, Lemma~\ref{lem:ext}, Lemma~\ref{lem:witt},
Lemma~\ref{lem:extshadow}, Theorem~\ref{thm:tbar},
Corollary~\ref{cor:profiles}, and the eliminations of
Sections~\ref{sec:p11}--\ref{sec:p03}.  The machine supplies exactly the
following finite facts, each re-derived from first principles at every
run: the Krawtchouk table \eqref{eq:kraw} and the row reductions behind
Proposition~\ref{prop:even}, Proposition~\ref{prop:id51} (whose
multipliers are stated there), and Proposition~\ref{prop:universal}; the
$6$-dimensional Lagrangian census ($1395$ subspaces, $135$ Lagrangians,
$15$-fold point incidence); the $O^+(8,2)$ shadow incidence of
Lemma~\ref{lem:extshadow} ($135$ Lagrangians through $r$, $16$ singular
shadow classes each, $72$ eligible classes of incidence exactly $30$, and
the $120$-element transvection orbit); the $90$ translation pairs and the
$4005$-pair compatibility exhaustion of Lemma~\ref{lem:p01-bins}; the
disjointness equivalence over all $9828$ weight-three words; the
$81{,}081$-candidate normal-form exhaustion of Lemma~\ref{lem:lc}; and the
four collision-degree histograms of Proposition~\ref{prop:p03}.

\subsection{The artifact}
The proof package contains five result-producing implementations, written
independently in different representations: a packed-symplectic-integer
checker ($33$ named checks), a literal-Pauli clean-room checker in Python,
an exact-rational-algebra checker built on
\texttt{fractions.Fraction}, a tuple-based finite-geometry enumerator,
and a separately authored JavaScript enumerator (Node.js).  A single
command,
\begin{center}
\texttt{python3 -I proof/replay.py}
\end{center}
runs all five from fresh processes, byte-compares their five JSON
certificates against re-runs, performs $35$ independent exact
cross-comparisons, runs $27$ Python proof tests---including $16$ primary
mutation fixtures, one per compared claims region, each of which must be
rejected---and the JavaScript clean-room mutation test (six internal
checks), executes two independently written
$\code{n,k,d}$ verifiers on calibration codes (the $\code{5,1,3}$ code and
a published $\code{14,3,4}$ code, whose distances $3$ and $4$ are
recomputed by full $2^{17}$-element enumeration), and writes an
SHA-256 manifest over all sources and outputs.  Everything is exact
integer and rational arithmetic; there is no floating point, no random
seed, no network access, and no external solver.  The artifact is
available at
\url{https://github.com/Mansehej/quantum-14-3-5}
and as ancillary files with this paper.

\subsection{Independent adversarial verification}
Beyond its internal cross-checks, the package was subjected to an
independent adversarial verification campaign (also AI-driven, with
implementations written from scratch without consulting the package
code).  That campaign re-derived the complete constraint algebra of
Sections~\ref{sec:even}--\ref{sec:profiles} from first principles;
re-enumerated all finite-geometry and collision counts (obtaining, e.g.,
exact maximum induced edge counts $12,12,13,27$ against the conservative
bounds $12,12,13,28$ used above); and, most stringently, validated the
distance-free lemmas empirically on random instances, where every object
can be enumerated directly: on ten random odd and seven random all-even
$11$-dimensional isotropic subspaces of $\F_2^{28}$, and on the real
$\code{14,3,4}$ calibration code, the shadow formula \eqref{eq:eshadow},
the $135$-extension self-duality and average \eqref{eq:nbar}, the
per-word $30/135$ incidence law, and the average \eqref{eq:tbar} were
confirmed exactly, word for word.  The campaign found no mathematical
flaw.  We stress the epistemic status honestly: all of this is machine
verification of machine-generated mathematics.  This paper is the
package's exposure to human review, and the finite certificates are
designed so that a skeptical reader can re-derive any of them in
isolation.

\section{Concluding remarks}\label{sec:remarks}

Theorem~\ref{thm:main} settles the entry $(n,k)=(14,3)$ of the stabilizer
code tables at $d=4$ exactly, resolving a question open since 1998
\cite{CRSS1998,Grassl}.  Two features of the proof seem worth isolating.
First, the two averaging identities---\eqref{eq:nbar} over the $135$
self-dual extensions and \eqref{eq:tbar} over their shadows---are
completely general for odd stabilizers with $k=3$ (and their analogues
exist for every $k$): they convert the existence of \emph{many} self-dual
codes above a given code into linear conditions on the code itself, at
the price of one finite incidence computation in an orthogonal geometry.
Second, the endgame is a purely combinatorial collision count among
weight-three shadow words, of a kind that seems applicable to other
small-parameter gaps where linear programming with integrality stalls.
A natural next step is formalization of the lemma layer in a proof
assistant; the finite certificates are already in a form amenable to
that.

\subsection*{Acknowledgements}
The computations described in Section~\ref{sec:verification} were
performed with open tooling on commodity hardware.  [Optional: add
acknowledgements.]

\begin{thebibliography}{9}

\bibitem{CRSS1998}
A.~R. Calderbank, E.~M. Rains, P.~W. Shor, and N.~J.~A. Sloane,
\emph{Quantum error correction via codes over GF(4)},
IEEE Trans.\ Inform.\ Theory \textbf{44} (1998), no.~4, 1369--1387.
arXiv:quant-ph/9608006.

\bibitem{Gottesman1997}
D.~Gottesman,
\emph{Stabilizer codes and quantum error correction},
Ph.D. thesis, California Institute of Technology, 1997.
arXiv:quant-ph/9705052.

\bibitem{Grassl}
M.~Grassl,
\emph{Bounds on the minimum distance of linear codes and quantum codes},
online tables, \url{http://www.codetables.de}, accessed 2026-07-26.

\bibitem{Rains1999Shadow}
E.~M. Rains,
\emph{Quantum shadow enumerators},
IEEE Trans.\ Inform.\ Theory \textbf{45} (1999), no.~7, 2361--2366.
arXiv:quant-ph/9611001.

\bibitem{RainsSloane1998}
E.~M. Rains and N.~J.~A. Sloane,
\emph{Self-dual codes},
in: V.~S. Pless and W.~C. Huffman (eds.), Handbook of Coding Theory,
North-Holland, Amsterdam, 1998, pp.~177--294.

\bibitem{LamPless1990}
C.~W.~H. Lam and V.~Pless,
\emph{There is no $(24,12,10)$ self-dual quaternary code},
IEEE Trans.\ Inform.\ Theory \textbf{36} (1990), no.~5, 1153--1156.

\bibitem{NRS2006}
G.~Nebe, E.~M. Rains, and N.~J.~A. Sloane,
\emph{Self-Dual Codes and Invariant Theory},
Algorithms and Computation in Mathematics 17, Springer, Berlin, 2006.

\bibitem{Grove2002}
L.~C. Grove,
\emph{Classical Groups and Geometric Algebra},
Graduate Studies in Mathematics 39, American Mathematical Society,
Providence, RI, 2002.

\bibitem{BallCentellesHuber}
S.~Ball, A.~Centelles, and F.~Huber,
\emph{Quantum error-correcting codes and their geometries},
Ann.\ Inst.\ Henri Poincar\'e D \textbf{10} (2023), no.~2, 337--405.
arXiv:2007.05992.

\end{thebibliography}

\end{document}
